%% Annexe C : Modèles et procédures
\chapter{Modèles et Procédures}
\label{ann:templates}

\section*{Fiche de scellé}
\begin{tableaucours}{L{3.4cm}L{8.8cm}}{%
    \tableauHeader{Rubrique} & \tableauHeader{Contenu}
}
Numéro de scellé & Identifiant unique \\
Pièce & Description, photographie \\
Saisie & Date, lieu, agent \\
Empreinte & Condensat à la réception \\
Transferts & De qui à qui, quand, pourquoi \\
\end{tableaucours}

\section*{Fiche d'acquisition (image)}
\begin{tableaucours}{L{3.4cm}L{8.8cm}}{%
    \tableauHeader{Rubrique} & \tableauHeader{Contenu}
}
Support & Référence, type \\
Méthode & Bloqueur, outil, image bit à bit \\
Empreintes & Condensat original et copie \\
Opérateur & Qui, quand \\
\end{tableaucours}

\section*{Fiche de constat}
\begin{tableaucours}{L{3.4cm}L{8.8cm}}{%
    \tableauHeader{Rubrique} & \tableauHeader{Contenu}
}
Constatation & Le fait établi \\
Méthode & Comment il a été établi, norme suivie \\
Preuve & Artefact, condensat, source \\
Interprétation & Analyse, niveau de confiance \\
\end{tableaucours}

\section*{Plan du rapport forensique}
\begin{enumerate}
    \item Synthèse pour le magistrat.
    \item Cadre et mission (réquisition, périmètre).
    \item Méthode (normes, outils, actes, condensats).
    \item Constatations, par domaine.
    \item Interprétation et niveaux de confiance.
    \item Limites et NEK.
    \item Annexes (chaîne de possession, condensats, journaux).
\end{enumerate}
